\documentclass[11pt]{letter}

\usepackage[margin=0.5in]{geometry}
\usepackage{hyperref}

\begin{document}

\begin{letter}{}

\opening{Dear Hiring Team,}

I am writing to apply for the Software Engineer, Machine Learning Platform, New Grad role at Quora. With a Master of Science in Computer Science from the University of Calgary and hands-on experience building Python/C++ systems for machine learning, large-scale data processing, performance-sensitive research software, and backend services, I am excited by the opportunity to work on infrastructure that helps ML engineers build, serve, and improve models at Quora's scale.

What draws me to this role is the mix of machine learning, distributed systems, developer velocity, and performance. My graduate research required building reliable software around complex data and models, not only training models in isolation. I developed Python and C++ pipelines for ingesting, preprocessing, storing, retrieving, and evaluating large Earth observation datasets, including DGGS-based workflows that convert imagery into structured tensor data. I also improved one C++ processing workflow from 233 seconds to 26 seconds for a 10-tensor workload through multithreading and performance-focused implementation.

I also bring practical ML experience from my DEM super-resolution work, where I trained PyTorch ResNet models, built data collection and preprocessing workflows using Google Earth Engine and Rasterio, and evaluated model quality with terrain-aware metrics such as RMSE, slope, TRI, TPI, and wavelet-based errors. I built supporting tools in Streamlit and Polyscope to inspect data and model outputs, which strengthened my interest in the surrounding systems that make ML development faster, more reliable, and easier to reason about.

My software background includes backend engineering as well. As a back-end developer, I built Django REST APIs for NLP, text-to-speech, and speech-to-text services backed by PostgreSQL, implemented authentication flows, and contributed to a WebRTC-based chatbot system using GPT-style responses and FAQ data. As a teaching assistant, I supported students in big data applications, data-at-scale workflows, SQL, AWS-based data pipelines, and database systems, which helped me communicate technical ideas clearly while debugging real implementation issues.

I have not yet worked professionally with Go, Kubernetes/EKS, NVIDIA Triton, Ray, or large-scale model serving infrastructure, but the role's note that previous ML infrastructure experience is not required is exactly why it feels like a strong match. I bring solid Python and C++ foundations, PyTorch experience, Docker and AWS exposure, strong learning habits, and a genuine interest in infrastructure that makes ML systems more scalable, reliable, and efficient. I am also comfortable with Quora's 9am to 3pm Pacific coordination hours from Calgary.

Thank you for considering my application. I would welcome the opportunity to discuss how my background in Python/C++ engineering, ML workflows, performance optimization, data pipelines, and research software can contribute to Quora's Machine Learning Platform team.

\closing{Sincerely,

Amir Mirzai Golpayegani

\href{mailto:amirgolpaa24@gmail.com}{amirgolpaa24@gmail.com}}

\end{letter}

\end{document}
